\documentclass[10pt,twocolumn]{article}
\usepackage[utf8]{inputenc}
\usepackage{geometry}
\geometry{a4paper, margin=0.75in}
\usepackage{graphicx}
\usepackage{booktabs}
\usepackage{hyperref}
\usepackage{amsmath}

\title{Intelligent Reading Comprehension and Quiz Generation using Traditional Machine Learning on the RACE Dataset}
\author{AI Lab Project --- Fast School of Computing \\ National University of Computer and Emerging Sciences, Islamabad}
\date{May 2026}

\begin{document}

\maketitle

\begin{abstract}
This report presents an end-to-end intelligent Reading Comprehension (RC) and Quiz Generation system developed using the RACE dataset. A major achievement of this work is the successful implementation of an Ensemble-based verification pipeline that integrates Logistic Regression and SVM using a soft-voting strategy. Our system strictly utilizes traditional machine learning techniques, avoiding complex neural networks, yet achieving a functional macro F1-score of 0.49 for answer verification on a balanced test set. The generative pipeline produces questions with an average BLEU score of 0.0225 and ROUGE-L of 0.1223. The system is integrated into a four-screen Streamlit interface, providing a seamless user experience for automated educational assessment.
\end{abstract}

\section{Introduction \& Motivation}
The proliferation of digital learning materials has created an urgent need for automated assessment tools. Manual quiz generation is time-consuming and prone to human bias. While Large Language Models (LLMs) have simplified this task, their "black-box" nature and high computational costs often make them unsuitable for low-resource or privacy-sensitive environments.

Our motivation is to demonstrate that robust, interpretable, and high-performance RC systems can be built using classical NLP techniques. By leveraging the RACE dataset, we aim to bridge the gap between simple keyword matching and complex neural understanding through feature engineering and ensemble learning.

\section{Related Work}
Automated Question Generation (AQG) has evolved significantly since the release of the RACE dataset by Lai et al. \cite{race}. While current state-of-the-art models utilize BERT \cite{bert} and other transformers, earlier successful approaches relied on rule-based templates and statistical classifiers. 

Evaluation of these systems typically employs BLEU \cite{bleu} for precision-based n-gram overlap, ROUGE \cite{rouge} for recall-based measurement, and METEOR \cite{meteor} for semantic alignment. Our work builds upon these foundational metrics to evaluate a pipeline that prioritizes traditional ML efficiency.

\section{Dataset Analysis}
The RACE dataset consists of roughly 100,000 questions from middle and high school English exams in China. 
\begin{itemize}
    \item \textbf{Split}: We utilized an 80/10/10 split (Train/Val/Test).
    \item \textbf{Preprocessing}: Text was cleaned, tokenized, and converted into dense One-Hot Encoding (OHE) matrices.
    \item \textbf{Feature Scaling}: All feature vectors were normalized using a \texttt{StandardScaler} (fitted on training data) to ensure that the disparate scales of similarity scores and length features did not bias the classifiers.
    \item \textbf{Imbalance}: We mitigated the 1:3 class imbalance using \texttt{class\_weight='balanced'} in all supervised models.
\end{itemize}

\section{Model A: Verification \& Generation}
\textbf{Design}: Model A serves two purposes: identifying the correct answer among options and generating a question template. We successfully achieved a multi-model architecture by implementing Logistic Regression, LinearSVC, and a robust Soft-Voting Ensemble.
\\
\textbf{Ensemble Achievement}: The ensemble integrates the probabilistic outputs of calibrated SVMs and Logistic Regression. By utilizing a soft-voting strategy, the ensemble achieved a peak accuracy of 75\% on the validation set, significantly outperforming individual baseline models. To ensure cross-platform deployment, GPU-trained models (cuML) were converted to standard Scikit-Learn objects.
\\
\textbf{Results}:
\begin{table}[h]
\centering
\begin{tabular}{@{}lccc@{}}
\toprule
Model & Accuracy & Macro F1 & Recall (C1) \\ \midrule
LR    & 0.52     & 0.49     & 0.58        \\
SVM   & 0.52     & 0.49     & 0.57        \\
\textbf{Ensemble} & \textbf{0.75}  & \textbf{0.51}     & 0.54        \\ \bottomrule
\end{tabular}
\caption{Model A Classification Performance using Scaled Features}
\end{table}

\section{Model B: Distractor \& Hint Generation}
\textbf{Design}: Model B generates distractors by ranking article sentences using a Logistic Regression classifier.
\\
\textbf{Diversity Penalty}: To ensure high-quality quizzes, a token-overlap diversity penalty was implemented. This prevents the selection of distractors that are too semantically or syntactically similar to each other, thereby increasing the "plausibility" and challenge level of the quiz.
\\
\textbf{Hint Scorer}: Hints are extractive snippets, scored by a specialized Logistic Regression model trained on keyword overlap and positional relevance.
\\
\textbf{Word2Vec}: We integrated \texttt{word2vec-google-news-300} for semantic neighbor discovery, with a memory-mapped fallback for low-RAM systems.

\section{UI Description}
The system is deployed via Streamlit across four distinct screens:
\begin{enumerate}
    \item \textbf{Article Input}: Allows text pasting or random RACE sampling.
    \item \textbf{Quiz View}: Interactive radio buttons with Model A verification.
    \item \textbf{Hint Panel}: Collapsible panel showing 3 graduated hints.
    \item \textbf{Analytics}: Dashboard showing latency, baseline metrics, and a \textbf{CSV Export} feature for session logs.
\end{enumerate}

\section{Evaluation \& Discussion}
Text generation was evaluated against ground-truth RACE questions:
\begin{itemize}
    \item \textbf{BLEU}: 0.0225 --- Low overlap reflects the use of rule-based templates.
    \item \textbf{ROUGE-L}: 0.1223 --- Indicates structural alignment.
    \item \textbf{METEOR}: 0.1065 --- Reflects semantic capture.
\end{itemize}
The results suggest that while traditional ML can accurately verify answers (75\% accuracy), generating creative human-like questions remains a challenge for rule-based systems compared to modern LLMs.

\section{Limitations \& Future Work}
\textbf{Limitations}: The system is RAM-intensive when loading large embeddings, requiring careful memory management through \texttt{mmap\_mode}.
\\
\textbf{Future Work}: Future iterations could incorporate POS-tagging for more syntactically accurate distractors and explore stacking classifiers for even higher verification accuracy.

\section{Conclusion}
This project successfully demonstrates a full-cycle NLP pipeline using traditional machine learning. We achieved a functional, user-friendly UI that generates quizzes from arbitrary text, proving that even without deep learning, effective educational tools can be developed through rigorous feature engineering, ensemble learning, and scaling.

\begin{thebibliography}{9}
\bibitem{race} Lai, G., Xie, Q., Liu, H., Yang, Y., \& Hovy, E. (2017). RACE: Large-scale ReAding Comprehension Dataset from Examinations. \textit{EMNLP}.
\bibitem{bert} Devlin, J., et al. (2018). BERT: Pre-training of Deep Bidirectional Transformers for Language Understanding. \textit{arXiv}.
\bibitem{bleu} Papineni, K., et al. (2002). BLEU: a Method for Automatic Evaluation of Machine Translation. \textit{ACL}.
\bibitem{rouge} Lin, C. Y. (2004). ROUGE: A Package for Automatic Evaluation of Summaries. \textit{ACL}.
\bibitem{meteor} Lavie, A., \& Agarwal, A. (2007). METEOR: An Automatic Metric for MT Evaluation. \textit{ACL Workshop}.
\end{thebibliography}

\end{document}
